% ========== EXTENSION EXECUTION MODELS ==========
% In-process vs out-of-process execution: The Pit vs The Grand Stage
% Source: Symphony/Content/The Pit, The Grand Stage, The In-Process, The Out-of-Process

\section{Extension Execution Models}
\label{sec:extension-execution-models}

\lettrine{S}{ymphony's} architecture recognizes that different types of extensions have fundamentally different performance and safety requirements. Rather than forcing all extensions into a single execution model, Symphony provides two distinct environments: \textbf{The Pit} for ultra-high-performance infrastructure operations, and \textbf{The Grand Stage} for safe, isolated user-facing functionality.

\subsection{The Dual Execution Philosophy}
\label{subsec:dual-execution-philosophy}

The choice between in-process and out-of-process execution represents one of the most critical architectural decisions in Symphony's design. This decision directly impacts performance, reliability, security, and developer experience.

\begin{infobox}[title=Execution Model Philosophy]
\textbf{Performance where it matters, safety where it counts}: Infrastructure operations that require nanosecond response times run in-process for maximum performance, while user-facing extensions run out-of-process for maximum safety and modularity.
\end{infobox}

This dual approach allows Symphony to achieve the best of both worlds: the raw performance of monolithic systems for critical operations, combined with the reliability and modularity of microkernel architectures for user-facing functionality.

\subsection{In-Process Extensions: The Pit}
\label{subsec:in-process-extensions-pit}

The Pit represents Symphony's high-performance execution environment where five specialized Rust extensions handle the infrastructure operations that form the foundation of intelligent orchestration.

\subsubsection{The Five Core Extensions}

The Pit consists of exactly five extensions, each with a specific responsibility in Symphony's infrastructure:

\begin{table}[h]
\centering
\caption{The Pit Extensions and Their Responsibilities}
\label{tab:pit-extensions}
\begin{tabular}{@{}lll@{}}
\toprule
\textbf{Extension} & \textbf{Primary Function} & \textbf{Performance Target} \\
\midrule
Pool Manager & AI model lifecycle management & 50μs allocation time \\
DAG Tracker & Workflow dependency tracking & 25μs state updates \\
Artifact Store & Data storage and retrieval & 15μs metadata operations \\
Arbitration Engine & Resource conflict resolution & 75μs decision making \\
Stale Manager & System cleanup and optimization & 10μs memory reclamation \\
\bottomrule
\end{tabular}
\end{table}

\textbf{Pool Manager: The Resource Maestro}

The Pool Manager handles every AI model in Symphony, from loading them into memory to shutting them down when complete. It implements sophisticated lifecycle management with predictive pre-warming and intelligent resource allocation.

Key capabilities include:
\begin{compactlist}
    \item \textbf{Predictive Loading}: Learns usage patterns and pre-warms models before requests
    \item \textbf{Health Monitoring}: Continuous model health checks with automatic failover
    \item \textbf{Cost Optimization}: Balances performance requirements with resource costs
    \item \textbf{Lifecycle Orchestration}: Manages the complete model lifecycle from idle to active
\end{compactlist}

\textbf{DAG Tracker: The Workflow Cartographer}

The DAG Tracker maintains real-time awareness of all workflow dependencies and execution states. It provides the coordination intelligence that enables Symphony's complex AI orchestration.

Core functions:
\begin{compactlist}
    \item \textbf{Dynamic Routing}: Adjusts execution paths based on real-time conditions
    \item \textbf{Critical Path Analysis}: Identifies and optimizes workflow bottlenecks
    \item \textbf{Recovery Planning}: Creates automatic checkpoints for failure recovery
    \item \textbf{Pattern Learning}: Optimizes future workflows based on historical execution
\end{compactlist}

\textbf{Artifact Store: The Institutional Memory}

The Artifact Store manages all data flowing through Symphony with intelligent versioning, quality scoring, and predictive caching.

Advanced features:
\begin{compactlist}
    \item \textbf{Quality Scoring}: Rates artifacts on structure, semantics, and utility
    \item \textbf{Predictive Preloading}: Anticipates data needs and preloads accordingly
    \item \textbf{Smart Versioning}: Efficient storage with complete change tracking
    \item \textbf{Relationship Mapping}: Tracks data lineage and dependencies
\end{compactlist}

\textbf{Arbitration Engine: The Wise Judge}

When multiple workflows compete for limited resources, the Arbitration Engine makes fair and intelligent allocation decisions using multi-dimensional optimization.

Decision-making capabilities:
\begin{compactlist}
    \item \textbf{Multi-Dimensional Analysis}: Considers business value, user impact, and system load
    \item \textbf{Fair Distribution}: Ensures equitable resource access over time
    \item \textbf{Adaptive Strategies}: Adjusts allocation policies based on system state
    \item \textbf{Learning from Outcomes}: Improves decisions based on historical results
\end{compactlist}

\textbf{Stale Manager: The System Curator}

The Stale Manager intelligently manages system resources and training data, prioritizing learning value over simple cleanup policies.

Curation strategies:
\begin{compactlist}
    \item \textbf{Training Value Assessment}: Evaluates artifacts for AI model improvement potential
    \item \textbf{Intelligent Retention}: Applies different policies based on data value
    \item \textbf{Storage Optimization}: Manages hot, warm, and cold storage tiers
    \item \textbf{Lifecycle Orchestration}: Handles complete data flow from creation to archival
\end{compactlist}

\subsubsection{In-Process Performance Characteristics}

The Pit's in-process architecture delivers exceptional performance through several key mechanisms:

\textbf{Direct Memory Access}

All Pit extensions share memory space with the Conductor Core, enabling zero-copy data transfers and eliminating serialization overhead.

\begin{table}[h]
\centering
\caption{In-Process Performance Metrics}
\label{tab:in-process-performance}
\begin{tabular}{@{}lll@{}}
\toprule
\textbf{Operation} & \textbf{Latency} & \textbf{Throughput} \\
\midrule
Model allocation & 50μs & 20,000/sec \\
Workflow state update & 25μs & 40,000/sec \\
Artifact metadata store & 15μs & 66,000/sec \\
Conflict resolution & 75μs & 13,000/sec \\
Memory reclamation & 10μs & 100,000/sec \\
\bottomrule
\end{tabular}
\end{table}

\textbf{Lock-Free Algorithms}

The Pit extensions use sophisticated lock-free data structures and atomic operations to achieve high concurrency without blocking.

\textbf{Cache-Optimized Memory Layout}

Data structures are carefully designed to maximize CPU cache efficiency, with cache-line alignment and memory access pattern optimization.

\subsubsection{Safety in Unsafe Territory}

Despite running in-process, The Pit maintains safety through several mechanisms:

\textbf{Rust Memory Safety}

All Pit extensions are written in Rust, providing compile-time memory safety guarantees that prevent common sources of crashes and security vulnerabilities.

\textbf{Panic Boundaries}

Rust panics are caught at the Python-Rust boundary, preventing extension failures from crashing the Conductor Core.

\textbf{Resource Isolation}

Even within the same process, extensions use separate resource pools and cannot interfere with each other's allocations.

\subsection{Out-of-Process Extensions: The Grand Stage}
\label{subsec:out-of-process-extensions-grand-stage}

The Grand Stage represents Symphony's realm of user-facing extensions, where safety, modularity, and developer experience take precedence over raw performance.

\subsubsection{The Three Extension Categories}

The Grand Stage hosts three types of extensions, each serving different aspects of the development experience:

\textbf{Instruments: AI \& ML Models}

Instruments bring intelligence to workflows through AI and ML models. Every AI model—whether hosted, local, or custom—becomes an Instrument extension.

Key characteristics:
\begin{compactlist}
    \item \textbf{Unified Interface}: All models use consistent input/output schemas
    \item \textbf{Configurable Intelligence}: Flexible parameter and cost limit settings
    \item \textbf{Maestro Mode}: Conductor-driven automatic model selection and blending
    \item \textbf{Manual Mode}: Complete user control over configuration and execution
\end{compactlist}

\textbf{Operators: Workflow Logic}

Operators handle data transformation, flow control, and utility functions. They provide the logical glue that connects different parts of workflows.

Capabilities include:
\begin{compactlist}
    \item \textbf{Data Preparation}: Cleaning, normalization, and structure transformation
    \item \textbf{File Conversion}: Seamless format transformation between different data types
    \item \textbf{Decision Logic}: Conditional flow control without custom scripting
    \item \textbf{Optimization}: Performance tuning and resource cleanup operations
\end{compactlist}

\textbf{Addons (Motifs): Visual Expression}

Addons create the user interface and experience layer, handling everything from custom dashboards to specialized editors.

Features:
\begin{compactlist}
    \item \textbf{Custom Dashboards}: Tailored data visualization and exploration interfaces
    \item \textbf{Specialized Editors}: Domain-specific editing environments
    \item \textbf{Developer Tools}: Debugging, profiling, and analysis utilities
    \item \textbf{UI Enhancements}: Theme customization and interaction improvements
\end{compactlist}

\subsubsection{Four-Layer Communication Architecture}

The Grand Stage implements a sophisticated four-layer architecture that bridges Python intelligence with Rust performance and React user interfaces:

\begin{figure}[h]
\centering
\begin{tikzpicture}[node distance=1.5cm]
    % Layer 1: Conductor Core
    \node[draw, rectangle, fill=brandPrimary!20, minimum width=8cm, minimum height=1.2cm] (conductor) {Conductor Core (Python RL Model)};
    
    % Layer 2: IPC Bus
    \node[draw, rectangle, fill=brandSecondary!20, minimum width=8cm, minimum height=1.2cm, below=of conductor] (ipc) {IPC Bus (Rust Communication Backbone)};
    
    % Layer 3: Extension Processes
    \node[draw, rectangle, fill=brandAccent!20, minimum width=8cm, minimum height=1.2cm, below=of ipc] (extensions) {Extension Processes (Isolated Execution)};
    
    % Layer 4: User Interface
    \node[draw, rectangle, fill=brandTertiary!20, minimum width=8cm, minimum height=1.2cm, below=of extensions] (ui) {User Interface (React + Shadcn)};
    
    % Arrows
    \draw[->] (conductor) -- (ipc) node[midway, right] {FFI Bridge};
    \draw[->] (ipc) -- (extensions) node[midway, right] {Unix Sockets};
    \draw[->] (extensions) -- (ui) node[midway, right] {Virtual DOM};
    \draw[<-] (conductor) -- (ipc);
    \draw[<-] (ipc) -- (extensions);
    \draw[<-] (extensions) -- (ui);
\end{tikzpicture}
\caption{Grand Stage Four-Layer Architecture}
\label{fig:grand-stage-architecture}
\end{figure}

\textbf{Layer 1: Conductor Core}

The Python-based reinforcement learning model that makes strategic orchestration decisions and manages global system state.

\textbf{Layer 2: IPC Bus}

A high-performance Rust crate that handles all inter-process communication with multiple transport optimizations:

\begin{table}[h]
\centering
\caption{IPC Transport Strategy Matrix}
\label{tab:ipc-transport-strategies}
\begin{tabular}{@{}llll@{}}
\toprule
\textbf{Message Type} & \textbf{Transport} & \textbf{Latency} & \textbf{Use Case} \\
\midrule
Control Messages & Unix Sockets & 0.1-0.3ms & Extension lifecycle \\
High-Frequency Data & Shared Memory & 0.01-0.05ms & Real-time UI updates \\
Large Payloads & Memory-mapped Files & 0.5-2ms & File transfers \\
Broadcast Events & Pub/Sub Multiplexing & 0.2-0.5ms & System notifications \\
\bottomrule
\end{tabular}
\end{table}

\textbf{Layer 3: Extension Processes}

Isolated processes that run individual extensions with resource limits and crash protection.

\textbf{Layer 4: User Interface}

React-based interface that renders Virtual DOM structures generated by Rust extensions, providing seamless integration between backend logic and frontend presentation.

\subsubsection{Process Isolation Benefits}

The out-of-process model provides several critical advantages:

\textbf{Crash Resilience}

Extension failures are completely contained within process boundaries. When an extension crashes, it can be automatically restarted without affecting other components.

\textbf{Resource Control}

Each extension runs with enforced memory and CPU limits, preventing resource exhaustion and ensuring predictable system behavior.

\textbf{Security Sandboxing}

Extensions run with minimal privileges and cannot access resources outside their granted permissions, providing strong security isolation.

\textbf{Hot Swapping}

Extensions can be updated, replaced, or reconfigured without system restart, enabling zero-downtime maintenance and deployment.

\subsection{Hybrid Model Strategy}
\label{subsec:hybrid-model-strategy}

Symphony's hybrid execution model represents a carefully balanced approach that optimizes for both performance and safety based on the specific requirements of different system components.

\subsubsection{Decision Criteria}

The choice between in-process and out-of-process execution follows clear criteria:

\textbf{In-Process (The Pit) Criteria:}
\begin{compactlist}
    \item Performance requirements below 100 microseconds
    \item Infrastructure operations that affect system-wide state
    \item Components written in memory-safe languages (Rust)
    \item Operations that require shared memory access
    \item Functions that cannot tolerate process restart latency
\end{compactlist}

\textbf{Out-of-Process (Grand Stage) Criteria:}
\begin{compactlist}
    \item User-facing functionality that can tolerate millisecond latencies
    \item Extensions developed by third parties or community
    \item Components that handle untrusted data or user input
    \item Functionality that benefits from independent scaling
    \item Operations that can be safely restarted without data loss
\end{compactlist}

\subsubsection{Performance Comparison}

The performance characteristics of each execution model reflect their different optimization priorities:

\begin{table}[h]
\centering
\caption{Execution Model Performance Comparison}
\label{tab:execution-model-comparison}
\begin{tabular}{@{}lll@{}}
\toprule
\textbf{Metric} & \textbf{The Pit (In-Process)} & \textbf{Grand Stage (Out-of-Process)} \\
\midrule
Function call latency & 100 nanoseconds & 0.1-0.3 milliseconds \\
Memory overhead & Shared with core & 50-200MB per extension \\
Startup time & Immediate & 100-500 milliseconds \\
Crash recovery & Process restart required & Automatic extension restart \\
Security isolation & Language-level & OS-level process isolation \\
Resource monitoring & Shared metrics & Per-process metrics \\
\bottomrule
\end{tabular}
\end{table}

\subsubsection{Migration Strategies}

Symphony's architecture allows components to migrate between execution models as requirements change:

\textbf{Pit to Grand Stage Migration}

Components can be moved from in-process to out-of-process execution when:
\begin{compactlist}
    \item Performance requirements become less stringent
    \item Security isolation becomes more important
    \item Third-party development is desired
    \item Independent scaling is beneficial
\end{compactlist}

\textbf{Grand Stage to Pit Migration}

Components can be moved from out-of-process to in-process execution when:
\begin{compactlist}
    \item Performance requirements become critical
    \item The component becomes stable and trusted
    \item Shared memory access is needed
    \item Process restart overhead becomes problematic
\end{compactlist}

This flexibility ensures that Symphony's architecture can evolve as the system matures and requirements change, always maintaining the optimal balance between performance and safety for each component.